% !TEX root = ../../../main.tex
\section{Direct sums and diagonalization criteria}
\label{sec:linear-algebra-i:direct-sums}

% Sources: IMG_9328--IMG_9330; October 22 handwritten PDF, pp. 1--3.
% The October 15 PDF does not contain this material.
\subsection{Transversal subspaces}

Suppose \(S,T\) are subspaces of \(V\) and \(S\cap T=\{0\}\).
If \(s\in S\), \(t\in T\), and \(s-t=0\), then
\(s=t\in S\cap T\), so both vectors are zero.

% !TEX root = ../../hw07.tex
% Source: ma-gy7033fall2026homework7 (1).pdf, p. 1, Problem 3.
% "Without relying on the lectures, prove" becomes "Prove directly"
% so that the shared statement is self-contained in the book.
\begin{exercise}
  \label{ex:hw07:disjoint-subspaces-independence}
  (10 points) Prove directly that if \(V_1,V_2\) are subspaces of \(V\) and
  \[
    V_1\cap V_2=\{0\},
  \]
  then if \(v_1\in V_1\) and \(v_2\in V_2\) are nonzero vectors,
  then \(\{v_1,v_2\}\) is a linearly independent set.
  \solutionlink{hw07:disjoint-subspaces-independence}
\end{exercise}


For a family \(S_1,\ldots,S_k\), write
\[
  S_1+\cdots+S_k
    =\{s_1+\cdots+s_k\mid s_j\in S_j\}.
\]
This is the span of the subspaces together.

\begin{definition}
  A family of subspaces \(S_1,\ldots,S_k\) is transversal if
  \[
    S_j\cap\sum_{i\ne j}S_i=\{0\}\qquad\text{for every }j.
  \]
\end{definition}
% Retain the course's defined use of transversal. Do not replace the
% family condition by pairwise trivial intersections, which is weaker.

The coordinate axes are transversal. In contrast, three distinct lines
through zero in a plane are not transversal, as illustrated in
\autoref{fig:linear-algebra-i:three-lines-not-transversal}.

\begin{marginfigure}[3\baselineskip]
  \centering
  % The October 22 p. 1 example: three distinct lines in the real plane.
\begin{tikzpicture}[
  x=10mm,y=10mm,font=\footnotesize,line width=0.45pt,
  continued/.style={{Latex[length=1.6mm,width=1.1mm]}-{Latex[length=1.6mm,width=1.1mm]}}
]
  \draw[continued,black!55] (-1.7,0) -- (1.7,0) node[right=2pt] {\(S_1\)};
  \draw[continued,black!55] (0,-1.5) -- (0,1.7) node[above=2pt] {\(S_2\)};
  \draw[continued] (-1.4,-1.05) -- (1.5,1.125) node[above right=2pt] {\(S_3\)};
  \fill (0,0) circle[radius=1.3pt];
  \node[below right=3pt] at (0,0) {\(0\)};
  \path (current bounding box.south) ++(0,-2mm);
\end{tikzpicture}

  \caption[Three lines and transversality]{A real-plane schematic.
    Each pair of distinct lines meets only at zero, but any two span the
    plane. Thus \(S_3\cap(S_1+S_2)=S_3\ne\{0\}\).}
  \label{fig:linear-algebra-i:three-lines-not-transversal}
\end{marginfigure}

\begin{proposition}
  \label{thm:linear-algebra-i:transversal-unique-components}
  Let \(S_1,\ldots,S_k\) be transversal, with \(s_j\in S_j\).
  If \(s_1+\cdots+s_k=0\), then every \(s_j=0\).
  Consequently, if also \(t_j\in S_j\) and
  \(s_1+\cdots+s_k=t_1+\cdots+t_k\), then \(s_j=t_j\) for every \(j\).
\end{proposition}

\begin{proof}
  Isolate the last component in the zero relation. It satisfies
  \[
    s_k=-\sum_{j=1}^{k-1}s_j
      \in S_k\cap(S_1+\cdots+S_{k-1})=\{0\}.
  \]
  Hence \(s_k=0\). Repeating with the remaining components gives
  \(s_1=\cdots=s_k=0\).

  For two equal sums, subtract them. Each \(s_j-t_j\) belongs to
  \(S_j\), and their sum is zero. The first conclusion gives
  \(s_j-t_j=0\) for every \(j\).
\end{proof}

\begin{definition}
  If \(S_1,\ldots,S_k\) are transversal and
  \(S_1+\cdots+S_k=V\), then \(V\) is their direct sum, written
  \[
    V=S_1\oplus\cdots\oplus S_k=\bigoplus_{j=1}^k S_j.
  \]
\end{definition}

\subsection{Combining bases}

For two complementary lines in \(\R^2\), each vector has one component
on each line, as in \autoref{fig:linear-algebra-i:direct-sum-components}.
\begin{marginfigure}
  \centering
  % Two-dimensional real schematic; the two subspace lines span the plane.
\begin{tikzpicture}[
  x=11mm,y=11mm,font=\footnotesize,line width=0.45pt,
  vector/.style={-{Latex[length=1.6mm,width=1.1mm]}},
  continued/.style={black!50,line width=0.35pt,{Latex[length=1.6mm,width=1.1mm]}-{Latex[length=1.6mm,width=1.1mm]}},
  guide/.style={dashed,black!45,line width=0.35pt}
]
  \coordinate (O) at (0,0);
  \coordinate (S) at (1.8,0.3);
  \coordinate (T) at (0.35,1.4);
  \coordinate (V) at ($(S)+(T)$);
  \draw[continued] (-0.6,-0.1) -- (2.7,0.45) node[right=2pt] {\(S\)};
  \draw[continued] (-0.15,-0.6) -- (0.6,2.4) node[above=2pt] {\(T\)};
  \draw[guide] (S) -- (V) -- (T);
  \draw[vector] (O) -- (S) node[midway,below=3pt] {\(s\)};
  \draw[vector] (O) -- (T) node[midway,left=3pt] {\(t\)};
  \draw[vector] (O) -- (V) node[above right=2pt] {\(v\)};
  \node[below left=2pt] at (O) {\(0\)};
  \path (current bounding box.south) ++(0,-2mm);
\end{tikzpicture}

  \caption[Components in a direct sum]{A real-plane schematic of
    \(V=S\oplus T\). The decomposition \(v=s+t\), with \(s\in S\)
    and \(t\in T\), is unique. The summands need not be perpendicular.}
  \label{fig:linear-algebra-i:direct-sum-components}
\end{marginfigure}

\begin{proposition}
  \label{thm:linear-algebra-i:direct-sum-bases}
  If \(V=S_1\oplus\cdots\oplus S_k\), concatenating a basis of each
  \(S_j\) gives a basis of \(V\). In particular,
  \[
    \dim V=\sum_{j=1}^k\dim S_j.
  \]
\end{proposition}

\begin{proof}
  Let \((e_{j,1},\ldots,e_{j,n_j})\) be a basis of \(S_j\).
  Suppose a linear combination of all these vectors is zero,
  \[
    \sum_{j=1}^k\sum_{\ell=1}^{n_j}a_{j,\ell}e_{j,\ell}=0.
  \]
  Group the terms into vectors
  \(s_j=\sum_{\ell=1}^{n_j}a_{j,\ell}e_{j,\ell}\in S_j\).
  Transversality gives \(s_1=\cdots=s_k=0\). Since each group is a
  linear combination of a basis of \(S_j\), all its coefficients
  \(a_{j,\ell}\) are zero. Thus the concatenated list is independent.

  The subspaces together span \(V\), and the basis vectors in each
  group span its subspace. Hence the concatenated list spans \(V\).
  It is therefore a basis, with \(n_1+\cdots+n_k\) vectors.
\end{proof}

\subsection{Transversality of distinct eigenspaces}

\begin{theorem}[Transversal eigenspaces]
  \label{thm:linear-algebra-i:distinct-eigenspaces-transversal}
  Let \(\lambda_1,\ldots,\lambda_k\) be distinct eigenvalues of
  \(L:V\to V\). Their eigenspaces
  \(E_{\lambda_1}(L),\ldots,E_{\lambda_k}(L)\) are transversal.
\end{theorem}

\begin{proof}
  We show that a relation
  \[
    v_1+\cdots+v_k=0,\qquad v_j\in E_{\lambda_j}(L),
  \]
  forces every \(v_j=0\). The argument proceeds by induction on the
  number of eigenspaces. For one eigenspace, the relation already says
  \(v_1=0\).

  Suppose the result holds for \(k-1\) distinct eigenspaces. Apply
  \(L\) to the relation and subtract \(\lambda_k\) times the original
  relation. Since \(L(v_j)=\lambda_jv_j\), this gives
  \[
    \sum_{j=1}^{k-1}(\lambda_j-\lambda_k)v_j=0.
  \]
  The \(k\)-th term has disappeared. Each remaining vector
  \((\lambda_j-\lambda_k)v_j\) still belongs to \(E_{\lambda_j}(L)\).
  By the induction hypothesis, all these vectors are zero. The scalars
  \(\lambda_j-\lambda_k\) are nonzero, so \(v_j=0\) for \(j<k\).
  The original relation now gives \(v_k=0\).

  A vector in one eigenspace that is a sum of vectors in the others
  produces such a zero relation by moving all terms to one side.
  The result just proved forces that vector to be zero. This is
  transversality.
\end{proof}
% October 22 p. 2 supplies this elimination argument and says to repeat it.
% The induction spells out that supplied iteration, replacing p. 1's
% explicitly abandoned attempt. Pairwise intersections alone are insufficient.

\subsection{Criteria for an eigenbasis}

\begin{theorem}[Spanning eigenspaces and diagonalization]
  \label{thm:linear-algebra-i:spanning-eigenspaces-diagonalizable}
  If \(\lambda_1,\ldots,\lambda_k\) are eigenvalues of \(L:V\to V\)
  and
  \[
    E_{\lambda_1}(L)+\cdots+E_{\lambda_k}(L)=V,
  \]
  then \(L\) is diagonalizable.
\end{theorem}

\begin{proof}
  Keep one copy of each distinct eigenspace in the sum. By
  \autoref{thm:linear-algebra-i:distinct-eigenspaces-transversal}, these
  eigenspaces are transversal, so the spanning hypothesis gives a
  direct-sum decomposition of \(V\). Choose a basis in each eigenspace.
  By \autoref{thm:linear-algebra-i:direct-sum-bases}, their concatenation
  is a basis of \(V\). Every vector in this basis is an eigenvector.
  Thus \(L\) is diagonalizable.
\end{proof}

The converse also holds. An endomorphism is diagonalizable exactly when
its eigenspaces span \(V\).
% October 22 p. 3 asserts the converse without writing a separate proof.

\begin{corollary}
  \label{thm:linear-algebra-i:distinct-eigenvalues-diagonalizable}
  If \(\dim V=n\) and \(L:V\to V\) has \(n\) distinct eigenvalues
  in \(\mathbb F\), then \(L\) is diagonalizable over \(\mathbb F\).
  In particular, a two-by-two matrix with two distinct eigenvalues in
  \(\mathbb F\) is diagonalizable over \(\mathbb F\).
\end{corollary}

The criterion does not require computing the dimensions of the individual
eigenspaces. When the eigenspaces do not span the whole space, there is
no eigenbasis. The nonzero nilpotent matrix in
\autoref{ex:linear-algebra-i:same-polynomial-not-similar} is an example.

\subsection{Approximation by diagonalizable matrices}

\begin{proposition}
  \label{thm:linear-algebra-i:density-of-diagonalizable-complex-matrices}
  The matrices diagonalizable over \(\C\) are dense in
  \(\mathrm M_{n\times n}(\C)\).
\end{proposition}

Here density means that every complex matrix can be approximated
arbitrarily closely, entry by entry, by matrices diagonalizable over
\(\C\).
% October 22 p. 3 describes density in every matrix neighbourhood.
% On October 28 Lucas clarified that this is a statement over C.
% No proof is supplied; do not identify density with a full-measure claim.
